\documentclass[UTF8,a4paper,11pt,fontset=fandol]{ctexart}
\usepackage[margin=22mm]{geometry}
\usepackage{amsmath,amssymb,booktabs,longtable,xcolor,xurl,hyperref}
\hypersetup{colorlinks=true,linkcolor=teal,urlcolor=teal,pdftitle={Phi3Lab4.0: Graph-driven recursive BPH integrals through four loops}}
\setlength{\parskip}{.4em}\setlength{\emergencystretch}{3em}
\title{从图自动构造 BPH 有限积分\\\large 一至四圈图集、Gaussian Taylor 代数与独立验证}
\author{Phi3Lab 4.0\quad MPhys 软件研究协作记录}
\date{2026 年 10 月 19 日}
\begin{document}\pagestyle{plain}\maketitle
\begin{abstract}
本稿记录软件自动化程度的实质扩展：由生成图的边、顶点和外腿标签，自动识别 UV 子图并展开 Bogoliubov 递归，选择相容整数动量路由，构造多变量 Taylor 系数，再完成 Gaussian 圈积分及整体 Schwinger 尺度积分。新入口不调用既有逐拓扑积分公式。一至四圈共698张图均已编译；全部图及精确权重与官方 feyngen 互校。两圈结果与既有独立推导的有限积分比较，三、四圈选定实例实际求积。primitive 图的 Hepp 子集证书提供参数空间和原动量绝对收敛界，递归余项证书、实际数值和求积误差分别标明。本项目尚未核验的 Z 部分不作前提。
\end{abstract}
\paragraph{使用顺序}
先沿第二节的递归与 Gaussian 构造推导一圈例子，再检查自动生成的两圈重叠图矩阵；之后用 primitive 的三、四圈实例理解解析证书，最后阅读实际积分及复跑记录。图集、表达式、积分结果及严格误差证明属于不同层次，本文分别记录。
\setcounter{tocdepth}{1}
\tableofcontents\clearpage
\section{范围、图集与生成算法}
六维、有共同正质量平方 $M>0$ 的实标量 $\phi^3$ 理论；固定外腿标签、截肢、连通且1PI。保留平行边和允许的自环。每圈测度为 $d^6k/\pi^3$，耦合及图权重单列。两圈高精度专用路线保留为独立基准，新引擎的生产计算不导入那些拓扑核。
\begin{center}\begin{tabular}{lrrr}\toprule
圈数 & 一点图 & 二点图 & 三点图\\\midrule
1&1&1&1\\2&1&2&7\\3&3&10&56\\4&12&64&540\\\bottomrule
\end{tabular}\end{center}
总数698。高圈枚举用度约束的邻接矩阵逐行生成、未处理等色顶点的轨道剪枝及精确染色细化与个体化标号；平行边阶乘和自环翻转计入完整对称因子。官方 feyngen 的逐图、精确权重字典用于独立比较，而非只比较图数。
\section{A graph-driven recursive BPH Gaussian compiler}
\label{sec:automatic-compiler}

The input is a connected, amputated, one-particle-irreducible multigraph
$G$ of massive scalar $\phi^3$ theory in six Euclidean dimensions.
Internal edges retain distinct indices, including parallel edges; external
half-edges retain their physical labels.  Let $L$, $I$ and $E$ denote its
loop, internal-edge and external-leg counts.  Trivalence and Euler's
identity give
\[
 I=3L+E-3,\qquad \Omega_G=6L-2I=6-2E.
\]
The present implementation accepts $1\le L\le4$ and $1\le E\le3$.
Its generated expressions use no topology-specific integration kernel.
This section establishes the regulated construction and its interior
parameter formula, and states separately the analytic and numerical
evidence needed to identify and integrate its unregulated limit.

\subsection{The primary construction is the Bogoliubov recursion}

For an embedded edge subgraph $\gamma$, all omitted incident half-edges
are separate boundary ports.  This convention matters when parallel
edges have the same endpoints.  The compiler enumerates proper,
connected, 1PI edge subsets with
\[
 \omega_\gamma=6L_\gamma-2I_\gamma\ge0.
\]
There is no replacement of an edge subset by an induced vertex subset.
If $P_\gamma$ lists all independent boundary momenta, the Taylor operator
is the full total-degree polynomial
\[
 t_\gamma^{\omega_\gamma} f(P_\gamma)
 =\sum_{r=0}^{\omega_\gamma}\frac1{r!}
   \left.\frac{d^r}{dz^r}f(zP_\gamma)\right|_{z=0}.
\]
The boundary variables can include other loop momenta of $G$; they are
not restricted to the physical external momenta.

For a nonempty spinney $S$ of mutually vertex-disjoint proper divergent
subgraphs, put
\begin{align*}
 \overline R_G I_G
 &=I_G+\sum_{S\ne\varnothing}
       \left(\prod_{\gamma\in S} C_\gamma\right)I_{G/S},\\
 C_\gamma&=-t_\gamma^{\omega_\gamma}\overline R_\gamma I_\gamma,
 \qquad
 R_G I_G=(1-t_G^{\Omega_G})\overline R_G I_G.
\end{align*}
The local polynomial carried by a two-point counterterm is retained in
full, including its mass and kinetic terms.  Contracting a subgraph in
this notation does not discard that polynomial.

The actual program expands this recursion by choosing disjoint maximal
subgraphs and then recursively preparing their proper subgraphs.  It
records a word $W$ with its subgraphs ordered from inner to outer and
coefficient $(-1)^{|W|}$.  The maximal members and the successive maximal
members inside them are unique for any resulting word.  There is
therefore no additional multiplicity attached to a word.  For the
finite input graph the resulting support can be compared with a
separately enumerated set of compatible forests.  This is an auxiliary
finite combinatorial check; the implementation does not invoke an
unverified general forest or $Z$ identity to construct its answer.

\subsection{A compatible tree gives exact integer momentum coordinates}

For a fixed recursive word, its selected subgraphs are nested or
vertex-disjoint.  A spanning tree can be constructed inside each minimal
selected subgraph, extended inside its parent, and finally extended to
$G$.  Equivalently, the implemented Kruskal ordering first processes
edges belonging to the smallest selected subgraph, breaking ties by the
original edge index.  For every selected $\gamma$,
$T\cap E_\gamma$ is then a spanning tree of its vertices.

Let the edges outside $T$ be the chords $h_1,\ldots,h_L$.  Assigning a
unit momentum to each chord and solving conservation along tree leaves
gives
\[
 k_e=\sum_{j=1}^L r_{ej}\ell_j
       +\sum_{a=1}^{E-1}r_{e,L+a}p_a,
 \qquad r_{ej}\in\mathbb Z,
 \qquad p_E=-\sum_{a<E}p_a.
\]
The reduced oriented incidence matrix of a tree has determinant $\pm1$.
Consequently this solve is integral, and different fundamental cycle
coordinates are related by an integer unimodular transformation,
possibly followed by a shift linear in the physical momenta.  Its loop
integration Jacobian is one.  An internal chord of $\gamma$ has its tree
path entirely in $\gamma$, so its fundamental cycle vanishes on every
edge outside $\gamma$.  The compiler checks both the number of these
chords, $L_\gamma$, and this absence of leakage.

At a positive common Schwinger cutoff, loop integration and Taylor
differentiation may be interchanged.  Translation invariance of those
Gaussian integrations also makes the integrated local Taylor polynomial
independent of the compatible routing used for its word.  This does
not assert that an external-momentum-dependent loop shift commutes
pointwise with an unintegrated Taylor operator.

\subsection{Scaled rows implement the boundary Taylor operators}

For each selected $\gamma$ introduce a formal scalar $t_\gamma$.  For
each edge $e$ and each coordinate $j$ of the routing, define
\[
 \chi_\gamma(e,j)=
 \begin{cases}
 1,&e\in E_\gamma\ \text{and }j\text{ is not an internal chord of }\gamma,\\
 0,&\text{otherwise}.
 \end{cases}
\]
Every physical coordinate is external to $\gamma$ in this definition.
The row used to build its Gaussian form is
\[
 c_{ej}(t)=r_{ej}\prod_{\gamma\in W}
                 t_\gamma^{\chi_\gamma(e,j)}.
\]
Thus a boundary loop momentum is scaled even though it is not a physical
external momentum.  In a nested pair an inner physical coefficient can
contain $t_{\rm inner}t_{\rm outer}$, an outer internal loop coefficient
inside the inner subgraph contains only $t_{\rm inner}$, and an inner
internal loop coefficient is unchanged.  These are precisely the
inside-out boundary scalings required by the recursion.

For $\omega_\gamma=0$ the program sets $t_\gamma=0$ directly.  For the
other selected subgraphs it works in the finite ring
\[
 \mathbb R[t_\gamma:\gamma\in W,\ \omega_\gamma>0]
     /(t_\gamma^{\omega_\gamma+1}:\gamma\in W),
\]
and finally sums all coefficients in the rectangular index set
$0\le k_\gamma\le\omega_\gamma$.  A coefficient in this ring already
means $\partial^k f(0)/\prod k_\gamma!$.  Multiplying it by another
factorial would be an error.  The rectangular indexing concerns the
different Taylor operators; each individual operator is still a
total-degree Taylor polynomial in all of its boundary momenta.

\subsection{Gaussian integration in the finite Taylor algebra}

Give every original edge its own positive Schwinger parameter $\alpha_e$.
No parameter is removed when a recursive word inserts a counterterm.
For that word set
\[
 Q_W(t)=\sum_{e=1}^I\alpha_e c_e(t)c_e(t)^T
       =\begin{pmatrix}A&B\\B^T&C\end{pmatrix},\qquad
 U_W=\det A,
\]
where $A$ is the $L\times L$ loop block.  With the physical Gram matrix
$g_{ab}=p_a\cdot p_b$, define
\[
 H_W=\sum_{a,b=1}^{E-1}g_{ab}
                  (C-B^TA^{-1}B)_{ab}.
\]
The constant loop block is positive definite at an interior parameter
point: the internal subgraph cycles and the quotient cycles still form
a complete basis after its selected boundary momenta have been set to
zero.  Gaussian elimination is therefore legitimate in the Taylor
ring.  Inverting a pivot uses a finite geometric expansion in its
nilpotent remainder; a logarithm uses the corresponding finite Taylor
series.  These algebraic operations implement derivatives of the
Gaussian integral, rather than numerical finite differences.

For each of the six spatial coordinates the same scalar loop matrix
occurs.  With the convention $d^6\ell/\pi^3$ per loop,
\[
 \int\prod_{j=1}^L\frac{d^6\ell_j}{\pi^3}
       e^{-\sum_e\alpha_e k_e(t)^2}
       =U_W(t)^{-3}e^{-H_W(t)}.
\]
There is no diagram symmetry weight, coupling, $(4\pi)$ factor, or
dimensional-regularization prefactor in this equation.

\subsection{The overall subtraction makes the scale integral explicit}

Write $\alpha_e=\rho x_e$, $\sum_e x_e=1$.  The Jacobian is
$\rho^{I-1}d\rho\,d^{I-1}x$, $U(\rho x)=\rho^LU(x)$, and
$H(\rho x)=\rho H(x)$.  Hence the remaining scale power is
\[
 \rho^{I-1-3L}=\rho^{E-4}.
\]
The simplex measure here is the usual coordinate Lebesgue measure,
whose volume is $1/(I-1)!$.  No extra simplex factorial is to be inserted.

Since $H$ is quadratic in all physical momenta, the overall zero-degree
subtraction for $E=3$ removes $1$ from $e^{-\rho H}$; the full
degree-two subtraction for $E=2$ removes $1-\rho H$.  For $M=m^2>0$,
\begin{align*}
 \Phi_3(H,M)
 &=\int_0^\infty\frac{e^{-M\rho}}\rho
                    (e^{-H\rho}-1)d\rho
 =-\log(1+H/M),\\
 \Phi_2(H,M)
 &=\int_0^\infty\frac{e^{-M\rho}}{\rho^2}
                    (e^{-H\rho}-1+\rho H)d\rho\\
 &=(M+H)\log(1+H/M)-H.
\end{align*}
For $E=1$ momentum conservation leaves no independent external momentum;
the full zero-momentum Taylor condition therefore gives $R_GI_G=0$.
This includes the chosen finite condition on the linear counterterm.

For a word define $\mathcal T_W$ to set its zero-order variables to zero
and sum the retained Taylor coefficients of its other variables.
The projective density emitted by the compiler is
\[
 \mathcal D_G(x;p,M)=
   \sum_{W\text{ from the recursion}}(-1)^{|W|}
    \mathcal T_W\bigl[U_W(t)^{-3}\Phi_E(H_W(t),M)\bigr].
\]
All words are evaluated at the same original $x$ and combined before
numerical integration.  In particular a separately divergent word is
not assigned an unregulated numerical integral.

\subsection{A two-loop overlapping graph extracted from the actual program}

The following matrices were extracted by
\path{tools/extract_automatic_overlap.py} from actual
\texttt{compile\_bph} output, retained in
\path{docs/automatic-compiler-evidence/generated-overlap.json}.  They are an illustration of the
algorithm, not a formula table used by its engine.  The edge order is
\[
 (0,2),(0,3),(1,2),(1,3),(2,3),
\]
with physical legs at vertices $0,1$.  Write its parameters as
$(a,b,c,d,e)$ and $p^2=s$.  Its two proper divergent triangles have edge
sets $\gamma_a=\{0,1,4\}$ and $\gamma_b=\{2,3,4\}$.
They overlap, so the recursive words are only
$\varnothing$, $\{\gamma_a\}$ and $\{\gamma_b\}$.

For the empty word, the program obtains the routing rows in the basis
$(\ell_1,\ell_2,p)$
\[
 \begin{pmatrix}1&1&1\\-1&-1&0\\-1&0&-1\\1&0&0\\0&1&0\end{pmatrix}.
\]
The full loop-plus-physical matrices after each selected zero-order
Taylor substitution are
\[
 Q_\varnothing=
 \begin{pmatrix}
 a+b+c+d&a+b&a+c\\
 a+b&a+b+e&a\\
 a+c&a&a+c
 \end{pmatrix},
\]
\[
 Q_{\gamma_a}=
 \begin{pmatrix}c+d&0&c\\0&a+b+e&0\\c&0&c\end{pmatrix},
 \qquad
 Q_{\gamma_b}=
 \begin{pmatrix}a+b&0&-a\\0&c+d+e&0\\-a&0&a\end{pmatrix}.
\]
The compatible tree and chord basis of the last word differ from those
of the first word; each transformation has unit Jacobian.  Put
\begin{align*}
 U&=(a+b)(c+d)+e(a+b+c+d),\\
 V&=ab(c+d)+cd(a+b)+e(a+c)(b+d).
\end{align*}
Direct determinants and Schur complements of these emitted matrices give
\begin{align*}
 U_\varnothing&=U,&H_\varnothing&=sV/U,\\
 U_{\gamma_a}&=(c+d)(a+b+e),&H_{\gamma_a}&=s\,cd/(c+d),\\
 U_{\gamma_b}&=(a+b)(c+d+e),&H_{\gamma_b}&=s\,ab/(a+b).
\end{align*}
The resulting density is the first $U^{-3}\Phi_2$ expression minus
the other two.  In particular no simultaneous
$t_{\gamma_a}t_{\gamma_b}$ term is introduced for the overlapping pair.

\subsection{A common regulator and the exact scope of the limit}

For regulated algebra use the same propagator on every original edge,
\[
 D_\eta(k)=\int_\eta^\infty e^{-\alpha(k^2+M)}d\alpha,
 \qquad \eta>0.
\]
Every recursive word retains the same original region
$\alpha_e\ge\eta$.  All Gaussian integrals and finitely many Taylor
derivatives are then absolutely integrable.  After projectivization the
common lower scale limit is
\[
 T(x,\eta)=\eta/\min_e x_e.
\]
Replacing the lower limits in the two definitions of $\Phi_E$ by this
same $T$ gives the regulated Gaussian expression exactly.  For fixed
interior $x$ the subtracted exponential is $O(\rho)$ for $E=3$ and
$O(\rho^2)$ for $E=2$, also after any of the finitely many Taylor
derivatives.  The mass controls infinity.  It follows that its
$\eta\downarrow0$ limit is the displayed projective density.

That interior conclusion alone does not justify taking the limit
through the entire simplex or through the original $6L$-dimensional
momentum integral.  To make both identifications one needs suitable
integrable majorants for the complete recursive sum, uniform in the
common regulator.  A future inductive proof must also retain control in
the internal subgraph-external momenta, not just at one fixed reference
scale.  The separate one- and two-loop family proofs provide this
bridge for their stated domains.  Graph-derived primitive certificates
provide additional bounds for their explicitly certified graphs.
Compiling all three- and four-loop
graphs, or testing a collection of UV rays, is not a proof of the
corresponding uniform bridge for every such graph.  The compiler's
formal expression, a pointwise numerical UV diagnostic, an integrated
estimate, and an analytic convergence certificate are distinct claims.

\subsection{Arithmetic cancellation is checked separately from integration}

The common-parameter sum can be much smaller than its individual words.
For the two-loop nested bubble with edges
\[
 (0,1),(0,3),(1,2),(2,3),(2,3)
\]
and $s=M=1$, an actual randomized-integration point was
\[
 \begin{aligned}
 \alpha={}&(0.6311634306896491,\ 0.016292580874392757,\ 0.35254392304634735,\\
 &\qquad6.228782654682385\!\times10^{-8},\
 3.101784279669877\!\times10^{-9}).
 \end{aligned}
\]
Its unguarded binary64 density was $+440112873472$, while independent
100-digit differentiation gave $-3680.2843093271767$.  A positive
importance weight of about $1.71\times10^{-8}$ turned the false density
into a spurious weighted value about $7527$, rather than the correct
$-6.29\times10^{-5}$.  This was an arithmetic error, not statistical
evidence for a large positive tail.  A signed local BPH density is allowed
and is not clipped to impose positivity.

The initial safeguard recomputed the complete sum whenever a
proper-subgraph parameter sum or its complement was below $10^{-4}$.
That face check is retained, but simultaneous quadratic subtractions
can fail outside its threshold. The current whole-coefficient relative
trigger and optional checked native executor are described in
Section~\ref{sec:precision-repair}; both retain the independent standard
library Decimal path. If the smallest positive face scale, also
including $1$, is $r$, the Decimal starting precision is
\[
 d=\min\{470,\max\{80,32+12\lceil\max(0,-\log_{10}r)\rceil\}\}.
\]
It then recomputes with $\min(500,d+30)$ digits, requiring
\[
 |v_{d'}-v_d|\le10^{-25}\max\{|v_{d'}|,|v_d|\},
 \qquad d'=\min(500,d+30),
\]
with an explicit acceptance when both values are zero, and increases
precision again if necessary. Failure to stabilize by
500 digits raises an error. The displayed failing point now agrees
with the independent reference.  The batch path tries
\texttt{numpy.longdouble} first, but records its actual mantissa width;
on the tested platform it has only 53 significant bits.  No extra
precision is inferred from its type name.

Increasing the arithmetic precision was not by itself sufficient.  A
second test exposed a small-$z$ series truncated at $n=29$ for
$\Phi_2$, with $z=H/M$, that had been reused even in high-precision
arithmetic.  For parameters
\[
 (0.1,0.4,0.5,4\times10^{-21},6\times10^{-21}),\quad s=M=1,
\]
independent 160-digit differentiation gave
\[
 116357167448485560.1524657840956468920609\ldots,
\]
whereas both 110- and 200-digit calculations with that fixed series
stabilized at
\[
 116357167448126952.4750081595388072343536\ldots.
\]
Their common bias escaped an arithmetic precision comparison.  The
omitted analytic tail, before multiplication by $U^{-3}$, is bounded by
$M|z|^{30}/[30\cdot29(1-|z|)]$ when $|z|<1$.  It can therefore be
amplified near a UV face.  The Decimal implementation now uses
context-dependent convergent summation for both small-$z$ functions:
\[
 \log(1+z)=\sum_{n\ge1}\frac{(-1)^{n+1}z^n}{n},\qquad
 \Phi_2=M\sum_{n\ge2}\frac{(-1)^nz^n}{n(n-1)}.
\]
It continues until adding the next term no longer changes the sum in
the active context, then performs the higher-precision recheck.  Both
failure cases are regression tests, rather than being hidden by a
wider numerical tolerance.

Decimal recomputation evaluates the accepted binary64 input parameters
and Gram entries; it does not invent extra input accuracy. The face
threshold, whole-coefficient condition trigger and precision comparisons
are practical arithmetic diagnostics,
not outward-rounded error intervals or a theorem that all cancellation
patterns have been detected.

\subsection{Numerical integration and current verification}

The optional generic integrator maps independently scrambled Sobol
points to Gamma variables and normalizes them to a symmetric Dirichlet
distribution on the original $I$-edge simplex.  For parameter $\beta$,
\[
 p_\beta(x)=\frac{\Gamma(I\beta)}{\Gamma(\beta)^I}
                   \prod_e x_e^{\beta-1},\qquad
 J=\mathop{\mathrm E}_{p_\beta}
                   [\mathcal D_G(x)/p_\beta(x)].
\]
The full word sum is formed before multiplication by this importance
weight and before averaging.  Standard errors use the variation of the
independently scrambled replicate means, not an independence assumption
for the points within one scramble.  The first half of each nested
Sobol sequence supplies a second power diagnostic.  These quantities,
maximum weighted samples, native precision checks and actual Decimal
fallback counts are recorded
separately; none is a rigorous total quadrature enclosure.

All 698 generated one- through four-loop graphs compile.  Their recursive
supports agree with the independent finite forest-support check.
Independent exact rational Gaussian determinants agree with
matrix-tree and spanning-two-forest polynomials for the tested original
words, including selected three- and four-loop primitive graphs.
Closed one-loop finite densities, physical-label and edge/vertex
permutations, scalar/batch agreement, mass dimension, zero subtraction
points, the two historical cancellation counterexamples and the new
simultaneous-face cases in Section~\ref{sec:precision-repair} are tested. Old
topology-specific entry points are patched to raise during an additional
test.  These checks substantiate the implemented compiler and its
recorded numerical results; they do not close an arbitrary-loop BPH
convergence induction or certify all three- and four-loop numerical
integrals.

\section{An exact graph-derived certificate for primitive graphs}

This section applies to a connected, one-particle-irreducible, bare trivalent
three-point graph in six Euclidean dimensions.  Let $L$ be its loop number,
$I=3L$ its number of internal edges, and $M=m^2>0$ the common squared mass.
Each loop measure is $d^6k/\pi^3$; couplings and symmetry factors are excluded.
No general forest identity is assumed.  In the primitive case the recursive
BPH definition has no proper counterterms and reduces directly to
$R_G=1-t_G^0$.

For every nonempty proper edge subset $A$, including disconnected subsets,
define its cycle rank and ultraviolet gap by
\begin{equation}
 h(A)=|A|-|V|+c(V,A),\qquad d(A)=|A|-3h(A).
 \label{eq:automatic-gap}
\end{equation}
Here isolated vertices are included in $c(V,A)$ and cancel from $h(A)$.
The certificate checks the exact integer condition $d(A)>0$ for all such
$A$.  It returns a request for a recursive-remainder certificate if this
condition fails; no claim about subdivergent graphs is inferred from the
primitive argument.

\subsection{Hepp sectors and exact subset dynamic programming}

Let $U$ be the first graph polynomial, with the convention
$U=\sum_T\prod_{e\notin T}\alpha_e$, summing over spanning trees.
Partition the positive projective domain by all edge orderings
$\alpha_{e_1}\leq\cdots\leq\alpha_{e_I}$.  In one such sector use the
projective gauge $\alpha_{e_I}=1$ and coordinates
\begin{equation}
 \alpha_{e_j}=\prod_{r=j}^{I-1}t_r,\qquad
 0<t_r<1,\qquad
 A_r=\{e_1,\ldots,e_r\}.
\end{equation}
The Jacobian is $\prod_{r=1}^{I-1}t_r^{r-1}$.  Ascending Kruskal selection
produces a spanning tree whose complement contains exactly $h(A_r)$ edges
from each prefix $A_r$.  Its positive monomial therefore gives
\begin{equation}
 U\geq\prod_{r=1}^{I-1}t_r^{h(A_r)},\qquad
 \int_{\text{sector}}\frac{\Omega}{U^3}
 \leq\prod_{r=1}^{I-1}\frac{1}{d(A_r)}.
 \label{eq:automatic-sector}
\end{equation}
There is no missing gauge factor: the projective form $\Omega/U^3$ has
degree $I-3L=0$, and its pullback to $\sum_e\alpha_e=1$ is precisely
$U^{-3}d\alpha_1\cdots d\alpha_{I-1}$.

The sum of the right-hand side of \eqref{eq:automatic-sector} over all
$I!$ orderings is computed without enumerating them.  With exact rational
arithmetic set
\begin{equation}
 D(\varnothing)=1,\qquad
 D(A)=\frac{\sum_{e\in A}D(A\setminus\{e\})}{d(A)}
 \quad(\varnothing\ne A\subsetneq E),\qquad
 B=\sum_{e\in E}D(E\setminus\{e\}).
 \label{eq:automatic-dp}
\end{equation}
Induction on $|A|$, separating an ordering by its last edge, proves this
recurrence.  Thus
\begin{equation}
 \int_{\sum\alpha=1}U^{-3}\,d^{I-1}\alpha\leq B<\infty.
 \label{eq:automatic-period-bound}
\end{equation}
The algorithm checks $2^I-2$ subsets and uses $O(I2^I)$ rational additions.
For the one-loop triangle, the primitive two-loop $K_{2,3}$ graph, the
three-loop cube with one vertex removed, and the four-loop Petersen graph
with one vertex removed, it gives respectively
\begin{equation}
 B=3,\quad24,\quad192,\quad804.
\end{equation}
These are rigorous upper bounds, not evaluations of the periods.  In
particular the previously derived exact two-loop period $1/2$ remains much
sharper than its graph-generic bound $24$.

\subsection{Finite BPH integral and uniform derivatives}

Use a fundamental-cycle routing and integrate its Gaussian loop variables.
Write $H$ for the nonnegative external Schur-complement quadratic form.
In the simplex gauge, a trial assignment setting all chord momenta to zero
gives
\begin{equation}
 0\leq H(p)\leq S,\qquad S=\max_{i=1,2,3}|p_i|^2,
 \qquad p_1+p_2+p_3=0.
 \label{eq:automatic-h-bound}
\end{equation}
Indeed the physical flow along each tree edge is a sum of a subset of
the three external momenta, hence is zero or one of $\pm p_i$.
The overall degree-zero Taylor subtraction and the common radial
Schwinger integration give
\begin{equation}
 J_G(p;M)=-\int_{\sum\alpha=1}U^{-3}
      \log\!\left(1+\frac{H(p)}{M}\right)d^{I-1}\alpha.
 \label{eq:automatic-primitive-bph}
\end{equation}
Consequently its absolute parameter integral, which here equals $-J_G$,
is bounded by $B\log(1+S/M)$.  This covers all simultaneous and
hierarchical ultraviolet faces, rather than only selected radial limits.

For a fixed physical shape set $J(\lambda)=J_G(\sqrt\lambda p;M)$,
$\lambda\geq0$.  For every $n\geq1$, differentiation under the integral is
justified by \eqref{eq:automatic-period-bound}, and gives
\begin{equation}
 |J^{(n)}(\lambda)|
 \leq B(n-1)!\left(\frac{S}{M+\lambda S}\right)^n.
 \label{eq:automatic-scale-derivative}
\end{equation}
These are derivatives of the exact integral; they are unrelated to a
quadrature-refinement error.  For a momentum direction $v$ with
$\sum_i v_i=0$ and $\max_i|v_i|\leq1$, positivity of the quadratic form
implies $H(v)\leq1$ and
$|H(p,v)|\leq\sqrt{H(p)H(v)}$.  At $\max_i|p_i|\leq P$ this yields
\begin{equation}
 |\partial_vJ_G|\leq\frac{2BP}{M},\qquad
 |\partial_u\partial_vJ_G|
 \leq B\left(\frac{2}{M}+\frac{4P^2}{M^2}\right)
 \label{eq:automatic-cartesian-derivatives}
\end{equation}
for two such directions $u,v$.  Replacing $M$ by a positive lower bound
$M_0$ makes all these estimates uniform on compact physical momenta and
all masses $M\geq M_0$.

\subsection{Original momentum absolute convergence and the regulator bridge}

Parameter convergence alone is not a proof of original momentum $L^1$
convergence.  The latter has a separate direct proof here.  Put
$D(Q)=(Q^2+M)^{-1}$ and
$I(k,p)=\prod_eD(Q_e(k,p))$.  In the above routing the physical coefficient
$C_ep$ vanishes on chords and is zero or $\pm p_i$ on tree edges.
The exact Taylor remainder is
\begin{equation}
 I(k,p)-I(k,0)=\int_0^1\frac{d}{dt}I(k,tp)\,dt,
\end{equation}
and each differentiated summand obeys
\begin{equation}
 \left|\frac{d}{dt}D(Q_e(k,tp))\right|
 \prod_{f\ne e}D(Q_f(k,tp))
 \leq 2|C_ep|\,D(Q_e(k,tp))^{3/2}
             \prod_{f\ne e}D(Q_f(k,tp)).
 \label{eq:automatic-fractional-majorant}
\end{equation}
This is a positive integrand.  Tonelli's theorem, the positive Schwinger
representation with one propagator power $3/2$, and Gaussian integration
give
\begin{align}
 &\int\prod_{\ell=1}^L\frac{d^6k_\ell}{\pi^3}
 D(Q_e)^{3/2}\prod_{f\ne e}D(Q_f)\\
 &\quad=\frac{\Gamma(1/2)}{\Gamma(3/2)}
 \int_{\sum\alpha=1}
 \frac{\alpha_e^{1/2}U^{-3}}{(M+H(tp))^{1/2}}\,d^{I-1}\alpha
 \leq\frac{2B}{\sqrt M}.
 \label{eq:automatic-fractional-integral}
\end{align}
The power of the common radial scale is $-1/2$, since
$\sum\nu_e=3L+1/2$.  Equations
\eqref{eq:automatic-fractional-majorant}--\eqref{eq:automatic-fractional-integral}
prove the explicit original-momentum estimate
\begin{equation}
 \|I(\cdot,p)-I(\cdot,0)\|_{L^1}
 \leq\frac{4B\sum_e|C_ep|}{\sqrt M}
 \leq\frac{4B(|V|-1)P}{\sqrt M}
 =\frac{8LBP}{\sqrt M}.
 \label{eq:automatic-momentum-l1}
\end{equation}

To justify the common original-edge regulator, replace every propagator by
$D_\eta(Q)=e^{-\eta(Q^2+M)}/(Q^2+M)
=\int_\eta^\infty e^{-\alpha(Q^2+M)}d\alpha$.
Writing $A=Q^2+M$, the differentiated factor contains
$(1+\eta A)e^{-\eta A}\leq1$, so exactly the majorant
\eqref{eq:automatic-fractional-majorant} works uniformly in $\eta>0$.
The joint $(t,k)$ majorant is integrable by
\eqref{eq:automatic-fractional-integral}; dominated convergence removes
the regulator on the original Taylor remainder.

For completeness, this domination can also be made common to all
$\max_i|p_i|\leq P$ and $M\geq M_0>0$: bounded tree shifts satisfy
$D(Q+a)\leq2(1+P^2/M_0)D(Q)$ for $|a|\leq P$.
Applying this inequality to the fractional majorants gives one integrable
majorant independent of the external momenta on the compact set.  If the
mass also varies over $M\geq M_0$, use the fixed unshifted propagator
$D_{M_0}$ on the right, since $D_M\leq D_{M_0}$.
On the parameter side, the regulated radial lower limit is
$\eta/\min_e\alpha_e$.  The absolute regulated logarithmic radial kernel
is bounded by $\log(1+H/M)\leq\log(1+P^2/M_0)$; hence
\eqref{eq:automatic-period-bound} supplies the second dominated-convergence
argument.  At finite $\eta$, Gaussian integration and Taylor subtraction
are legitimate.  The two limits therefore identify
\eqref{eq:automatic-primitive-bph} with the ordinary, absolutely convergent
momentum BPH remainder, in the stated normalization and subtraction
scheme.  No separately divergent unregulated term is integrated.

The software records the rational $B$ exactly.  Decimal bound values are
display approximations, not interval enclosures.  Neither these analytic
certificates nor numerical refinement comparisons constitute a rigorous
total quadrature-error estimate.

\clearpage\section{多 UV 面的算术修复与可复核执行路径}
\label{sec:precision-repair}

从图构造出正确的递归表达式之后，还必须可靠地计算其完整系数和。
本节记录四圈实际反例带来的修复。Bogoliubov BPH 递归、完整零动量
Taylor 条件及共同原始边参数均保持不变；新增机制处理的是有限表达式的算术，
不另选物理方案，也不调用逐拓扑预写的生产积分公式。

\subsection{为什么单面阈值遗漏了实际四圈错误}

\texttt{L4E3G205} 有三个顶点不交的二点 bubble 子图，原始边集分别为
$\{0,1\}$、$\{3,4\}$、$\{6,7\}$。令每个子图中的两条边各为
$\delta/2$，其余六条边各为 $(1-3\delta)/6$，并取
$M=p_1^2=p_2^2=p_3^2=1$。下列点都位于参数单纯形内部，
各子图参数和及补集参数和均大于原来的 $10^{-4}$ 阈值：
\begin{center}\small
\begin{tabular}{rll}\toprule
$\delta$ & 绕过保护的 binary64 密度 & 80/110 位参考及修复值\\\midrule
$10^{-2}$ & $-99552$ & $-99369.9951744098$\\
$10^{-3}$ & $-412316860416$ & $-81064258.2517440$\\\bottomrule
\end{tabular}
\end{center}
第一组相对误差约为 $0.183\%$，第二组相对误差约为 $5085$ 倍。
多个二阶 Taylor 减除同时作用时，共同抵消可以远强于每个单独面尺度所暗示的程度。
加大采样数不能消除这种系统算术误差。
不受保护的计算只用于重现历史缺陷，不作为被接受的积分密度。

\subsection{完整 Taylor 系数和的相对精度触发}

记第 $W$ 个递归项中保留的 Taylor 系数为
\[
 a_{W,\mathbf k}(\alpha)
 =[t^{\mathbf k}]\,U_W(t)^{-3}\Phi_E(H_W(t),M).
\]
除完整密度外，程序同时计算所有保留系数的绝对值之和：
\[
 \mathcal K_G=\sum_W(-1)^{|W|}\sum_{\mathbf k}a_{W,\mathbf k},
 \qquad
 A_G=\sum_W\sum_{\mathbf k}|a_{W,\mathbf k}|.
\]
原有 UV 面检查继续保留。对未被单面检查选中的点，若
\begin{equation}
 64\,\varepsilon_{\rm work}\,A_G
 >10^{-10}|\mathcal K_G|,
 \label{eq:relative-cancellation-trigger}
\end{equation}
或浮点计算得到非有限结果，就以同一输入重新计算完整递归和。
这里的绝对触发目标为零；$\varepsilon_{\rm work}$ 是实际工作类型的机器精度。
批量路径记录 \texttt{numpy.longdouble} 的真实尾数位数，
本机仍只有 53 个有效二进制位，不能从类型名称推定额外精度。

取同一 $\delta=10^{-3}$ 参数点，但令 $p_i^2=10^{-28}$，
未保护的值约为 $-4.1633\times10^{-17}$，而 110 位参考为
$-8.7361\times10^{-21}$。
若使用 $10^{-12}$ 的固定绝对阈值，该错误会因数值很小而被放过。
式~\eqref{eq:relative-cancellation-trigger} 会触发这个点，并恢复其相对精度。

$A_G$ 本身由有限精度算术求得，Jet 卷积和 Schur 消元的内部误差
没有被转换为严格的向外舍入区间。因此该规则是保守的数值触发条件，
并不是已经证明的全域向前误差界。
实际回归、未触发点的复查与独立解析参考继续分别记录。

\subsection{标准库复算与动态工作精度}

标准库 Decimal 保留完整且独立的执行路径。
把 $1$ 也计入各真子图及其补集参数和，记其中最小正数为 $r$，初始位数取
\[
 d=\min\{470,\max\{80,
           32+12\lceil\max(0,-\log_{10}r)\rceil\}\}.
\]
在该位数计算 $v_d$ 后，再以 $d'=\min(500,d+30)$ 位计算，要求
\[
 |v_{d'}-v_d|
 \le10^{-25}\max\{|v_{d'}|,|v_d|\}.
\]
两次均为精确零时接受；其余情况使用相对复核。
若不稳定则继续提高位数，最高 500 位仍未稳定时明确报错。
小参数的 $\log(1+z)$ 与 $\Phi_2$ 继续按当前上下文动态求和，
保留上一节两个真实历史反例的回归，避免固定项数级数在不同位数下共享同一偏差。
高精度算术使用已接收到的 binary64 输入及 Gram 数据，
不会增加原始输入本身的有效位数。

\subsection{可选原生执行器以 50/100 位逐点复核}

Python 从图生成子图、递归项和相容路由；可选 C++17 执行器只读取
生成程序中的整数路由系数、缩放单项式、Taylor 阶数及递归符号。
它不按图号选择拓扑公式。
Boost.Multiprecision 分别以 50 位和 100 位重新执行同一程序。
对非零结果，其接受检查包括
\[
 \max\left\{
  \frac{|v_{100}-v_{50}|}{|v_{100}|},\,
  \frac{10^{-90}A_{100}}{|v_{100}|}
 \right\}\le10^{-25}.
\]
零值另作明确检查。原生精度检查失败时，交给自适应 Decimal，
不静默采用未经确认的低精度值。
这同样是数值检查，不是严格舍入误差证明。

库加载器核对 ABI、操作系统与 CPU、当前 C++ 源码的 SHA256 和库文件的 SHA256；
缓存句柄在后续调用也重新核查。
库缺失、过期或属于不同主机时，标准库路径仍可工作。
\texttt{PHI3LAB\_DISABLE\_NATIVE=1} 可明确禁用原生后端。
该库由本项目随包源码构建，与独立参考所用的官方 pySecDec 库分别记录。

多精度结果输出到 binary64 时可能下溢，状态单独报告；
100 位文本接口保留未转为 binary64 的数值。
不能对 binary64 无法表示的非零结果承诺该输出类型的相对精度。

\subsection{实际验证和执行成本}

本次冻结源码下实际完成 15 个共同面、不等面、边和顶点重标号、
真嵌套及极小不变量回归，并另外检查 20 个实际未触发点。
80/110 位 Decimal、标量、批量和原生 50/100 位检查全部通过；
最大标量相对差约为 $2.17\times10^{-13}$，批量约为 $1.83\times10^{-13}$。
完整输入、程序哈希、库身份及触发计数保存在
\path{results/automatic_multiface_comparison.json}。

完整标准库套件在禁用 site packages 且设置
\texttt{PHI3LAB\_DISABLE\_NATIVE=1} 后实际执行 170 项，
明确跳过 34 项可选测试，全部执行项通过。
可选 NumPy、SciPy 和原生环境另实际执行 85 项，无跳过、失败或错误。
两个套件含重叠测试，计数不相加。
另外明确禁用原生后端的代表检查实际通过 8 项，跳过 3 项 NumPy 批量测试。
各份记录包含源码指纹及日志哈希。

1024 点实测基准包含 32 个共同 UV 面点和 992 个 Dirichlet $\beta=0.35$ 点。
110 位 Python Decimal 实际耗时 18.9802 秒；
同一生成程序的原生 1、2、4 线程分别为 3.5847、1.8712、0.99724 秒。
4 线程约快 19.03 倍，最大相对差为 $2.09\times10^{-16}$。
这些数字描述指定机器、输入和工作位数下的算术成本，
不能直接换算成所有完整积分的加速比。
原始每点记录和时间保存在 \path{docs/automatic-native-benchmark.json}；
冻结后的数学依赖哈希核对没有被称为重新计时。

本节的点态修复、完整振幅求积和解析收敛证书仍分别记录。
新增高精度执行路径使暴露的问题被实际重算和回归检查，
但点态数值触发、随机组误差及确定性求积细化都不因此成为严格的总积分误差界。

\section{Graph-derived importance sampling without factorial sectors}

The primitive certificate also yields an executable importance distribution.
This distribution is constructed from edge subsets of the input graph; it
does not select an integration formula by topology name. For every nonempty
proper subset put
\[
 d_+(A)=\max(1,|A|-3h(A)),\quad
 D_+(\varnothing)=1,\quad
 D_+(A)=\frac{\sum_{e\in A}D_+(A\setminus e)}{d_+(A)},\quad
 B_+=\sum_{e\in E}D_+(E\setminus e).
\]
Choose the largest edge with probability $D_+(E\setminus e)/B_+$, then
recursively choose the last edge of the remaining subset with probability
$D_+(A\setminus e)/\sum_fD_+(A\setminus f)$. This samples all edge orderings
without materializing $I!$ sectors. For its proper prefixes $A_r$ independently
draw $t_r=u_r^{1/d_+(A_r)}$. Set
\[
 \widetilde\alpha_{e_j}=\prod_{r=j}^{I-1}t_r,\quad
 \widetilde\alpha_{e_I}=1,\quad
 T=\sum_e\widetilde\alpha_e,\quad\alpha_e=\widetilde\alpha_e/T.
\]
The ordering probability telescopes to
$[B_+\prod_r d_+(A_r)]^{-1}$. Multiplying the conditional densities of the
$t_r$ cancels the factors $d_+$, whereas the simplex Jacobian is
$T^{-I}\prod_rt_r^{r-1}$. Therefore the complete, already combined BPH
density $\mathcal K_G$ has the unbiased weighted sample
\[
 B_+ T^{-I}\!\left(\prod_{r=1}^{I-1}t_r^{r-d_+(A_r)}\right)
 \mathcal K_G(\alpha).
\]
For primitive three-point graphs $d_+=d$ and $I=3L$. The proof in the
preceding section gives the pointwise envelope
\[
 |\text{weighted sample}|\le B\log(1+S/M).
\]
For graphs with subdivergences the clipped gaps specify a valid sampling
distribution only. They do not prove a bound on the combined recursive
remainder. The default implementation uses Hepp sampling for primitive
three-point inputs and Dirichlet sampling otherwise; both remain selectable
for comparison.

The independent alternative draws $z_e$ with a Gamma distribution of shape
$\beta$ and sets $\alpha_e=z_e/\sum z$. Its inverse density in the simplex
measure is
\[
 \frac{\Gamma(\beta)^I}{\Gamma(I\beta)}
 \prod_e\alpha_e^{1-\beta}.
\]
Independent scrambled Sobol replicates provide numerical means, between-run
standard errors and a nested half-power diagnostic. They are not rigorous
quadrature enclosures. On this platform NumPy's long double has only 53
effective binary bits; the explicit Decimal boundary recomputation is what
repairs the demonstrated cancellation failures.

The CDF implementation also has a recorded endpoint regression: every CDF
entry after the last active edge must equal one, not merely its final array
entry. Otherwise rounding can select an absent edge. Tests independently
check the normalized-simplex Jacobian, volume, first moments and primitive
weighted envelope, including the largest open-cube floating input.

\section{从图推导的分层重要性采样}
\label{sec:grouped-sampling}

This section changes the sampling distribution of the complete finite
recursive BPH density, not the renormalization prescription or its kernel.
The optional tool \texttt{tools/integrate\_bph\_groups.py} evaluates
\texttt{compile\_bph(graph).batch\_density}.  It imports neither a specialized
four-loop formula nor the independent momentum-space oracle.  The recursive
Bogoliubov subtraction, including every finite quadratic bubble term, is
therefore the same one used by the graph-driven compiler.  No unverified
forest identity is a premise of the construction below.

\subsection{Recognized structure and the change of variables}

The tool first checks every proper divergent edge subgraph.  Each must be a
two-edge, one-loop, two-point bubble of degree $\omega=2$, and their vertex
sets must be pairwise disjoint.  After contraction of these bubbles, the
bare complement must be a connected, one-particle-irreducible, one-loop
quotient.  These conditions are tested from the graph and its compiled
proper subgraphs, without selecting a formula by graph identifier.  Graphs
outside this class are explicitly rejected.  The class includes appropriate
two-, three-, and four-loop examples with two or three external legs.

Let $b$ be the number of bubbles, $M$ the number of bare complement edges,
and $N=2b+M$ the original number of edges.  Introduce bubble totals $r_i$,
a single bare total $A$, bubble splits $x_i$, and bare ratios $v_j$:
\begin{align*}
 &r_i>0,\quad A>0,\quad \sum_{i=1}^b r_i+A=1,\\
 &\alpha_{i,0}=r_i x_i,\qquad
   \alpha_{i,1}=r_i(1-x_i),\qquad 0<x_i<1,\\
 &\alpha_j=A v_j,\qquad v_j>0,\quad \sum_{j=1}^M v_j=1.
\end{align*}
Here the $\alpha_j$ in the last line refer only to bare complement edges.
In the original projective simplex Lebesgue measure, the absolute
Jacobian is
\begin{equation}
 d^{N-1}\alpha=
 \left(\prod_{i=1}^b r_i\right) A^{M-1}
 d^b(r_1,\ldots,r_b,A)\,
 \prod_{i=1}^b dx_i\,d^{M-1}v.
 \label{eq:hierarchical-jacobian}
\end{equation}
This follows by scaling the one independent coordinate of each two-edge
simplex by $r_i$ and the $M-1$ independent bare simplex coordinates by $A$.
It also follows directly by inserting the simplex delta functions before
performing the group scalings.

Write $G=b+1$ and sample independently
\[
 (r_1,\ldots,r_b,A)\sim\operatorname{Dirichlet}(a,\ldots,a),
 \qquad v\sim\operatorname{Dirichlet}(d,\ldots,d),
 \qquad x_i\sim\operatorname{Beta}(c,c).
\]
Dividing \eqref{eq:hierarchical-jacobian} by these probability densities gives
\begin{align}
 w={}&\frac{\Gamma(a)^G}{\Gamma(Ga)}
       \frac{\Gamma(d)^M}{\Gamma(Md)}
       \left[\frac{\Gamma(c)^2}{\Gamma(2c)}\right]^b
       \left(\prod_i r_i^{2-a}\right) A^{M-a}
       \left(\prod_j v_j^{1-d}\right)
       \prod_i[x_i(1-x_i)]^{1-c}.
 \label{eq:hierarchical-general-weight}
\end{align}
The default $a=d=1$, $c=3$ therefore has the particularly simple weight
\begin{equation}
 \boxed{\displaystyle
 w=\frac{\left(\prod_i r_i\right)A^{M-1}}
 {\Gamma(b+1)\Gamma(M)\,30^b
  \prod_i x_i^2(1-x_i)^2}.}
 \label{eq:hierarchical-weight}
\end{equation}
The implementation uses inverse Gamma transforms followed by normalization
for the two Dirichlet levels, and inverse Beta transforms for the splits.
There are $G+M+b=N+1$ Sobol coordinates.  The redundant coordinates in the
normalized Gamma constructions do not change their induced distributions.
For every compact interior integration domain the change of variables is
exact; for an absolutely integrable finite density $\rho_G$ it gives
\[
 \int_{\Delta_{N-1}}\rho_G(\alpha)\,d^{N-1}\alpha
       =\mathbb E\,[w\rho_G(\alpha)].
\]
This sampling identity is not a new general all-loop finiteness theorem.

\subsection{Why the bare total must be sampled separately}

At fixed other positive parameters, exact integration of a bubble's
internal momentum produces the quadratic BPH remainder
\[
 r^{-3}e^{-m^2r}
 \{e^{-r x(1-x)Q^2}-1+r x(1-x)Q^2\}.
\]
The identity
\[
 e^{-z}-1+z=z^2\int_0^1(1-t)e^{-tz}\,dt
\]
shows the factor $x^2(1-x)^2$.  At fixed external $Q$ it also gives the
local $r^{-1}$ behavior as $r\to0$, which is compensated by the pair
Jacobian $r$.  Other disjoint subtractions and the overall external Taylor
operation preserve this fixed-parameter split factor.  These observations
motivate $c=3$ and a flat density in each group total; they do not by
themselves prove a global finite-variance bound at simultaneous UV faces.

There is a second, different face.  All bare complement edges may shrink
together, $A\to0$, while bubble totals remain positive.  For a cubic graph
with one-loop quotient and $E$ external legs, $M=b+E$.  The term containing
all quadratic bubble counterterms has a contracted outer Gaussian factor
of order $A^{-3}$, and its polynomial loop moments can contribute
$A^{-b}$.  The overall finite radial kernels start with $H$ for $E=3$
and $H^2$ for $E=2$, with $H=O(A)$.  The resulting leading behavior can be
\[
 A^{-3-b}A^{4-E}=A^{-(M-1)}.
\]
The bare Jacobian in \eqref{eq:hierarchical-jacobian} then leaves an
integrable contribution of order $A^0$.  Treating the $M$ bare edges as
separate flat groups instead gives their sum a proposal density vanishing
as $A^{M-1}$.  Such a proposal badly undersamples this important face.
Ordinary edge-wise Dirichlet proposals can have the same problem: their
bare sum has density proportional to $A^{M\beta-1}$ near zero.  Sampling
$A$ as one top-level total keeps its default marginal density nonzero at
zero.  The failed flat-group pilot is retained as evidence rather than
discarded or described as successful variance control.

\subsection{Weight checks and numerical evidence}

An independent exact simplex moment is
\begin{equation}
 \int_{\Delta_{N-1}}\prod_e\alpha_e^{n_e}\,d^{N-1}\alpha
 =\frac{\prod_e\Gamma(n_e+1)}{\Gamma(N+\sum_e n_e)}.
 \label{eq:hierarchical-monomial}
\end{equation}
After the group substitution, the top-level powers are
$k_i=n_{i,0}+n_{i,1}+1$ and
$k_A=M-1+\sum_{\mathrm{bare}}n_j$.  The split Beta integrals and bare
simplex moment cancel the corresponding top-level Gamma numerators,
recovering \eqref{eq:hierarchical-monomial}.  The tool verifies pair
exponents $(2,2)$ and $(3,2)$, with zero bare exponents, by factorized
one-dimensional Gauss--Jacobi integrals at orders $32$ and $64$.  For the
default shapes the transformed factors are polynomials, and the observed
residual is assessed as floating-point rounding and refinement.  It also
checks the actual sampler and inverse weight at $256$ points against
independent Dirichlet and Beta probability-density implementations.

A constant-function importance Monte Carlo test is unsuitable for the
default split proposal: $w\sim x^{-2}$ and its second moment against
$x^2dx$ diverges.  The chosen monomials avoid that endpoint problem.  The
four-loop raw pilots, main and independent-seed runs, the separate
100-digit endpoint probe, and the derivation are retained in
\path{docs/automatic-group-sampling-evidence/}; final portable evidence is linked
from the accompanying release documentation.  Every run records the
actual tool hash, frozen core fingerprint, exact graph, compiled-program
hash, seeds, replicate values, half-power differences, and precision
fallback diagnostics.  The independent momentum-space scalar is used
only after integration for comparison.  Reported randomized errors and
refinement differences are numerical diagnostics, not rigorous total
quadrature enclosures.

\subsection{Recorded four-loop test of the proposal}

For the selected graph $\mathrm{L4E3G205}$, there are $b=3$ bubble groups
and $M=6$ bare edges.  Pilot variance and refinement diagnostics led to the
choice $a=\tfrac12$, $d=1$, $c=3$ in
\eqref{eq:hierarchical-general-weight}.  Its weight is
\[
 w=\frac{\pi^2}{\Gamma(6)\,30^3}
       \frac{(r_1r_2r_3)^{3/2}A^{11/2}}
       {\prod_{i=1}^3x_i^2(1-x_i)^2}.
\]
The kernel is still the unaltered generic recursive BPH sum.  At unit mass
squared and symmetric unit external invariants, the final catalog-aligned
tool was actually executed with $12$ scrambles of $2^{14}$ points and an
independent seed with $8$ scrambles of $2^{13}$ points:
\[
\begin{array}{c|r|c|c}
 &\text{value}&\text{reported standard error}
             &\text{half-power difference}\\ \hline
 \text{main}&-0.0005154702382511922&1.80923\,10^{-6}&1.18853\,10^{-6}\\
 \text{repeat}&-0.0005116648680529154&3.69832\,10^{-6}&2.66601\,10^{-7}
\end{array}
\]
The separately derived numerical momentum-space reference is
$-0.0005158785412715456$.  The absolute gaps are
$4.08303\,10^{-7}$ and $4.21367\,10^{-6}$, respectively, both below the
stated numerical target $10^{-5}$.  The main and repeat differ by
$3.80537\,10^{-6}$.  These comparisons validate this selected point
numerically; they do not turn the reported randomized errors into rigorous
intervals or extend the result to all higher-loop graphs.

The primary raw files are
\texttt{final-main-p14-r12.json} and
\texttt{final-repeat-p13-r8.json} in the portable evidence directory.
They were re-executed after the graph-numbering correction, rather than
assigning new hashes to earlier samples.  The final tool captures its
source hash at the start and verifies it again at the end, in addition to
checking the frozen production-core fingerprint.  Its generated-graph
indices agree with the public catalog; the earlier experimental ordering
and its two-loop calibration mapping are explicitly retained with the
historical evidence.  An independent audit checks the inverse density,
the full simplex Jacobian, monomial moments and graph relabelling.  A
separate $16$-case saved-record harness has one authentic positive control
and $15$ individual-corruption negative controls.  It accepts the positive
control and rejects all negative controls without executing any new integral.

\clearpage\section{四圈结果的独立外圈动量参考}
\label{sec:four-loop-independent}

This section concerns the fixed graph L4E3G205, with common squared mass
$M=1$ and symmetric Euclidean invariants $p_1^2=p_2^2=p_3^2=1$,
$p_1+p_2+p_3=0$.  Every original loop measure is $d^6k/\pi^3$; couplings,
vertex signs and diagram symmetry weights are excluded.  The reduction is
used only as an independent verification oracle.  Its implementation imports
neither the production Gaussian density and jets nor an existing
topology-specific integration kernel.  It does not replace graph-driven
production execution by a formula registry.

\subsection{The actual graph and the eight recursive terms}

The graph has nine vertices, twelve internal edges and external legs at
vertices $6,7,8$.  In original edge order its edges are
\begin{equation}
\begin{split}
 &(0,5),(0,5),(0,8),(1,4),(1,4),(1,7),\\
 &(2,3),(2,3),(2,7),(3,6),(4,5),(6,8).
\end{split}
\end{equation}
Its only proper divergent subgraphs are the three vertex-disjoint bubbles
on edge sets $\{0,1\}$, $\{3,4\}$ and $\{6,7\}$.  Directly from the
disjoint-spinney Bogoliubov recursion,
\begin{equation}
 \overline R_G I=
 (1-t_{\gamma_1}^{2})(1-t_{\gamma_2}^{2})(1-t_{\gamma_3}^{2})I,
 \qquad R_G I=(1-t_G^0)\overline R_G I.
 \label{eq:four-independent-eight}
\end{equation}
The product has eight terms.  Each local Taylor operator acts on the full
boundary momentum of its own bubble.  This is an elementary factorization
of this particular recursive expression, not an invocation of a general
forest or unverified project $Z$ identity.

Tracing the outer cycle gives
\begin{equation}
 6\longrightarrow\operatorname{bubble}(3,2)\longrightarrow7
 \longrightarrow\operatorname{bubble}(1,4)
 \longrightarrow\operatorname{bubble}(5,0)\longrightarrow8
 \longrightarrow6.
\end{equation}
The three side insertion counts are therefore $(1,2,0)$, rather than one
bubble on every side.  A side containing $n$ bubbles has $n+1$ ordinary
outer propagators, so define
\begin{align}
 \Sigma_R(s;M)
 &=\int_0^1\bigl[(M+s u(1-u))
       \log(1+s u(1-u)/M)-s u(1-u)\bigr],du,\\
 L_n(s;M)&=\frac{\Sigma_R(s;M)^n}{(s+M)^{n+1}}.
 \label{eq:four-independent-line}
\end{align}
The finite subtraction integrates to $-s/6$ and is retained throughout.
There is no extra factor $1/2$ for an insertion.  At finite common
original-edge Schwinger cutoff $\eta>0$, the independent inner bubble
integrations factor exactly in \eqref{eq:four-independent-eight}.

For the symmetric point we may use side counts $(0,1,2)$ and choose
$q_0=k$, $q_1=k+p_1$, $q_2=k-p_3$.  The complete overall degree-zero
subtraction is then
\begin{equation}
 J=\int\frac{d^6k}{\pi^3}
 \left[L_0(k^2)L_1((k+p_1)^2)L_2((k-p_3)^2)
       -\frac{\Sigma_R(k^2)^3}{(1+k^2)^6}\right].
 \label{eq:four-independent-outer}
\end{equation}
Writing three copies of $L_1$ would describe a different graph, although its
zero-point term would happen to be the same.  This makes the off-shell
branch mapping an essential part of the independent check.

\subsection{The three-dimensional integral and its stable evaluation}

Set $r=k^2$ and embed the external plane using
$p_1=(1,0)$ and $p_3=(-1/2,\sqrt3/2)$.  The projection of a uniform $S^5$
direction onto this plane has density
$2(1-x^2-y^2)/\pi$ on the unit disk.  With
$x=\sqrt v\cos\theta$, $y=\sqrt v\sin\theta$, the normalized measure and
shifted squared momenta are
\begin{align}
 \frac{d^6k}{\pi^3}
 &\longrightarrow\frac12 r^2,dr\,2(1-v)\,dv\,
                     \frac{d\theta}{2\pi},\\
 q_1^2&=r+1+2\sqrt r\,x,\qquad
 q_2^2=r+1+\sqrt r\,x-\sqrt{3r}\,y.
 \label{eq:four-independent-measure}
\end{align}
The directions $(x,y)$ and $(-x,-y)$ are paired before radial integration.
An exact independent closed form at $M=1$ is
\begin{equation}
 \Sigma_R(s)=\frac{(s+4)^{3/2}}{3\sqrt s}
       \operatorname{arsinh}(\sqrt s/2)-\frac43-\frac{4s}{9}.
 \label{eq:four-independent-bubble-closed}
\end{equation}
Near $s=0$, the implementation instead uses
\begin{equation}
 \Sigma_R(s)=\sum_{j=2}^{\infty}a_js^j,\qquad
 a_j=\frac{(-1)^j(j!)^2}{j(j-1)(2j+1)!},\quad
 a_2=\frac1{60},\quad
 \frac{a_{j+1}}{a_j}=-\frac{j-1}{2(2j+3)}.
\end{equation}
This closed form was also compared with a separate 70-digit
one-dimensional parameter integration at six scales; the largest observed
relative difference was $4.22\times10^{-15}$.

Let $\delta_\pm$ be the log ratio of the shifted product to its common
zero-point product, and set
$e=(\delta_++\delta_-)/2$, $o=(\delta_+-\delta_-)/2$.  The paired relative
difference is evaluated as
\begin{equation}
 \operatorname{expm1}(e)
       +e^e,2\sinh^2(o/2).
 \label{eq:four-independent-evenodd}
\end{equation}
At high radius the even and odd quantities are formed directly by analytic
Taylor derivatives and even/odd binomial grouping.  They are not obtained
by subtracting nearly equal floating logarithms.  The derivative
coefficients use 60--80-digit mpmath arithmetic before conversion to the
floating angular and radial quadrature; the full integral is not an
80-digit computation.

\subsection{A separate analytic radial-tail estimate}

For $r\geq4$, $|z|\leq1$ and $|a|=1$, write
$s=|k+za|^2$.  Then $r/4\leq s\leq9r/4$ and, with $\ell=\log r$,
the positive one-loop parameter formula gives
\begin{equation}
 0\leq\Sigma_R(s)\leq\frac{s\ell}{6},\qquad
 0\leq\Sigma_R'(s)\leq\frac\ell6,\qquad
 0\leq\Sigma_R''(s)\leq\frac1{6s}.
\end{equation}
Consequently, for $n=0,1,2$,
\begin{align}
 |L_n(s)|&\leq\frac{\ell^n}{6^n s},\\
 |L_n'(s)|&\leq\frac{(2n+1)\ell^n}{6^n s^2},\\
 |L_n''(s)|&\leq\frac{(4n^2+5n+2)\ell^n}{6^n s^3}.
 \label{eq:four-independent-line-bounds}
\end{align}
Put $g(z)=L_0(r)L_1(|k+zp_1|^2)L_2(|k-zp_3|^2)$.
Since $|s'|\leq3\sqrt r$ and $s''=2$, direct Leibniz differentiation gives
\begin{equation}
 |g''(z)|\leq\frac{20000}{27}\frac{\log^3r}{r^4}.
\end{equation}
The paired second-order Taylor remainder contributes a factor $1/2$;
the radial measure contributes another $1/2$.  Thus the whole omitted
outer tail above $T\geq4$ is bounded by
\begin{equation}
 |\operatorname{tail}(T)|\leq\frac{5000}{27T}
 \left(\log^3T+3\log^2T+6\log T+6\right).
 \label{eq:four-independent-tail}
\end{equation}
At $T=\operatorname{expm1}(48)$ this is
$3.109\times10^{-14}$, as a display approximation to the analytic bound.

Without pairing, the first Taylor derivative satisfies
$|g'|\leq(64/9)\log^3r/r^{7/2}$, which is also radially integrable.
For the common-cutoff bubble, its value and first two $s$ derivatives are
bounded by the same positive uncut quantities.  The cutoff outer
propagators have the corresponding uncut derivative majorants through
order two.  The above bounds therefore justify angular pairing and
dominated removal of $\eta$ in the reduced outer integral, identifying
\eqref{eq:four-independent-outer} with the same BPH regulated limit.
This argument is not, by itself, a general $L^1$ proof for the unreduced
24-dimensional combined momentum expression.

\subsection{A separate exact-coefficient Taylor truncation estimate}

Write $\Sigma_R(s)=s^2\phi(s)$ with
\begin{equation}
 \phi(s)=\int_0^1h(u)^2,du\int_0^1
             \frac{1-z}{1+s h(u)z},dz,\qquad h(u)=u(1-u).
\end{equation}
For $|w|=\varrho<1$, the ratio
$\phi(r(1+w))/\phi(r)$ is a positive-weight average of $(1+vw)^{-1}$,
$0\leq v\leq1$.  Its modulus is at most $(1-\varrho)^{-1}$ and its real
part is at least $(1-\varrho)/(1+\varrho)^2$.  It has a consistent analytic
logarithm, and
\begin{align}
 \left|\log\frac{\phi(r(1+w))}{\phi(r)}\right|
 &\leq2\log(1+\varrho)-\log(1-\varrho)+\frac\pi2,\\
 |\log(1+vw)|&\leq-\log(1-\varrho)+\frac\pi2.
\end{align}
Apply the two univariate Cauchy estimates separately to the shifted lines.
For counts $(0,1,2)$ there are three copies of $\log\Sigma_R$ and five
copies of $\log(1+s)$ in their combined log ratio.  At
$\varrho=9/10$ the sum of the boundary constants is less than $60$.

For a high-radius join $R$, Taylor order $J$, and
$q_0=(R^{-1}+2R^{-1/2})/\varrho<1$, put
$a=(J+1)/2$ and $D_*=60q_0/(1-q_0)$.  The Cauchy remainder, the exponential
mean-value estimate, and the radial prefactor bound
$\log^3r/432$ give
\begin{equation}
\begin{split}
 E_{\mathrm{Taylor}}\leq{}
 &\frac{e^{D_*}60q_0^{J+1}}{432(1-q_0)}\\
 &\times\left[
 \frac{\log^3R}{a}+\frac{3\log^2R}{a^2}
       +\frac{6\log R}{a^3}+\frac6{a^4}\right].
\end{split}
 \label{eq:four-independent-taylor-tail}
\end{equation}
For $R=1000,J=14$, this is $4.470\times10^{-15}$.
It bounds finite Taylor truncation with exact derivative coefficients.
It does not bound numerical coefficient rounding or the remaining
quadrature, and is distinct from \eqref{eq:four-independent-tail}.

\subsection{Calibration, actual refinements and preserved provenance}

The identical angular/radial scheme first reproduces the one-loop triangle
and the two-loop single-bubble BPH amplitude.  The latter's unit-point
reference is $-0.00617095361238479$.  In the last run its observed gap is
$1.27\times10^{-15}$; the one-loop calibration gap is
$6.25\times10^{-16}$.  The target four-loop QMC mean was not used to tune
the mapping, formula or refinement choices.

\begin{center}
\begin{tabular}{c c c}
\hline
Run & Radial/disk/half-angle & Four-loop value\\
\hline
Coarse & $12/20/24$ & $-0.0005158785412719462$\\
Fine & $20/32/40$ & $-0.0005158785412708389$\\
Finer & $28/48/64$ & $-0.0005158785412715456$\\
\hline
\end{tabular}
\end{center}
The first two runs use logarithmic panel width $2$, with respectively
$x=\log(1+r)\leq40,44$, join $R=10000$, Taylor orders $8,10$, and
60/70-digit derivative coefficients.  The last run uses panel width $1$,
$x\leq48$, join $R=1000$, Taylor order $14$ and 80-digit coefficients.
The last two results differ by $7.07\times10^{-16}$.  These are refinement
and calibration diagnostics, not a rigorous total quadrature enclosure.

The independent original script, three run JSON files, summary and
derivation notes are copied byte-for-byte into the documentation's
\texttt{automatic-four-loop-independent-evidence} directory, with original
times and hashes retained.  The portable tool
\path{tools/compare_automatic_four_loop_independent.py}
records its own hash and verifies the original script hash.  Its default
output is \path{results/automatic-four-loop-oracle-run.json}.
This reference is not bound to a transient production-source fingerprint;
a production comparison must preserve its own actual source, program,
execution and uncertainty records separately.

\clearpage\section{实际验证、证据范围与复现}

前面几节给出了从图到递归有限密度、解析证书和采样分布的连续构造。
本节说明哪些步骤已经实际执行，以及每种核查支持什么结论。
项目采用完整的有质量零动量 Taylor 条件，BPH 的 Bogoliubov 递归始终是主线；
尚未核验的一般 $Z$ 部分不作为数值结果或证明的前提。

\subsection{698 张图的覆盖意味着什么}

一至四圈、外腿数一至三的固定标签图共 698 张。全部图都已产生可执行表达式，
并在相等正参数点求值；每张图仍保存原始边索引，平行边可以分别进入子图。
全部规范图与精确对称权重也已经逐项匹配官方 feyngen，而不只匹配总数。
一、两圈保留旧图号；高圈图号采用新的精确染色规范序。
其他生成器中的图号只有在核对带固定外腿的同构关系后才能对应。

这里的完整覆盖是\emph{图枚举和表达式编译覆盖}。
698 个内部参数点都能求值，不等于 698 个完整振幅都已有高精度积分，
更不等于所有 UV 面附近的误差和原动量收敛都已统一证明。
这些层次在目录、执行记录和解析证书中分别保存。

\subsection{44 项独立有限密度检查}

独立参考程序从约化关联矩阵求路由，不复用生产程序的树叶消元；
它以 mpmath 构造 Gaussian 矩阵，再对边界缩放变量直接高精度微分，
不调用生产 Taylor jet、矩阵消元或密度算术。
它保留混合二次系数，包括多个嵌套二点反项的乘积。
另用精确生成树与生成二森林多项式检查原始 Gaussian 项的 $U,F$，
范围包含选定三圈、四圈 primitive 图。

实际复跑的 44 项检查全部通过，覆盖一、两圈完整图集的多个参数点，
以及高圈的选定 primitive、嵌套二点反项和 UV 条件控制实例。
这里的真嵌套由各检查图自身边序下的严格包含直接确认：三圈检查图的
$\{0,1\}\subsetneq\{0,1,2,3,8\}$；四圈检查图的
$\{0,1\}$ 与 $\{3,4\}$ 分别严格包含于 $\{0,1,2,3,4,5,6,7\}$。
它们分别对应新目录的 \texttt{L3E3G53} 与 \texttt{L4E3G300}，
各检查四种参数点模式；这与下述完整积分选用的图分开记录。
其中固定级数截断反例另用 160、200 位直接微分复核，
采用独立的 $10^{-80}$ 绝对阈值，避免一般相对容差掩盖大密度中的偏差。
对该点，生产 Decimal200 与高精度参考的差约为
$3.53\times10^{-122}$。
这些检查证实列出的代数与数值条件；它们不是整个 simplex 的积分误差包络。

\subsection{14 项完整两圈积分互校}

联合 UV 算术修复后，通用编译器已在当前冻结源码下重新实际执行
14 个两圈图与运动学输入的完整积分，
再与之前独立推导并核验的有限 BPH 公式比较。
专用公式只在验证工具中作为参考调用；通用生产计算仍由输入图自动生成。
全部 14 项比较通过，实际最大绝对差为
\[
 \max |J_{\rm generic}-J_{\rm reference}|
 =6.009245532560825\times10^{-6}.
\]
两圈验证工具先从完整UV子图判断分层提案是否适用。适用图使用固定总尺度形状$1/2$、
内部Beta形状3和bare边比例形状1，从12组各$2^{14}$点开始，
只有自身标准误不大于$5\times10^{-6}$且半样本差不大于$10^{-5}$才停止加密；
其余图保持原19次幂、12组的通用提案。参考值在选阶后才调用。
实际逐次加密历史与采样总量保存在结果中。此前低效普通提案的中断运行单独保留，
没有计作完成14项，也没有从屏幕打印伪造原始采样记录。
参考方案、每圈测度、质量平方、固定外腿和有限 Taylor 条件保持一致。
这一步比较的是完整有限积分，而不只是极点、$U,F$ 或某个内部参数点。
该最大差描述这 14 次运行；不能把它推广为所有运动学输入的严格统一精度。
独立随机组标准误、半样本细化差与实际参考差均保留在结果 JSON 中。

\subsection{35 张高圈 primitive 图的解析证书}

三圈 4 张、四圈 31 张 primitive 三点图，共 35 张，
都已实际构造图导出的 Hepp 子集证书。
证书检查全部真边子集的正 UV gap，并以精确有理动态规划计算 Hepp 系数。
上一节证明把该系数连接到参数绝对可积上界、原始动量 $L^1$ 控制，
以及共同 $D_\eta$ 正则化的移除；质量和外动量的适用范围随证书明确记录。

这 35 张是\emph{高圈 primitive 三点图}的解析覆盖，不能解释为所有高圈图已闭合。
对于有真发散子图的三圈、四圈图，裸图周期界不能直接控制完整递归余项。
采样时将非正 gap 截成正值只是定义重要性分布，不会自动成为递归 UV 证明。
一般的递归 UV 余项估计及其对内部子图外动量的统一控制仍有待完成。

\subsection{一次实际原生三圈有限 BPH 互校}

官方 pySecDec 1.6.6 已为 \texttt{L3E3G5} 实际生成并编译 C++ 原生库。
该图是新三圈三点目录的第 5 张图，包含 7 个顶点、9 条内部边，
物理外腿标签 $1,2,3$ 分别接在顶点 $4,5,6$，没有真发散子图，
其 Hepp 系数为 192。上游生成 300 个 sector；生成耗时约 32.99 秒，
原生编译耗时约 193.37 秒。随后分别在共同 $M=1$ 的零点和单位对称点，
以不同种子实际运行原生 QMC，耗时约 107.23 秒和 103.11 秒。

这里 $I=9,L=3,d=6-2\epsilon$。含完整 Gamma 前因子的参数表达式为
\[
 I_{\rm DR}(p,M)=\Gamma(3\epsilon)
       \int_\Delta U^{-3+4\epsilon}F(p,M)^{-3\epsilon}.
\]
令 $P=\int_\Delta U^{-3}$，其共同极点是 $P/(3\epsilon)$。
若 $c_0$ 表示完整有限系数，则
\[
 c_0(p,M)=\int_\Delta U^{-3}
       \left[\frac43\log U-\log F(p,M)-\gamma_E\right].
\]
由于图中没有真发散子图，BPH 递归只剩整体零阶减除。
在 $M=1$、$F(0,1)=U$ 下，严格有限转换正是
\[
 J_R(p,1)=c_0(p,1)-c_0(0,1)
 =-\int_\Delta U^{-3}\log\frac{F(p,1)}U.
\]
没有额外的 3 或 $1/3$ 因子；也没有乘入图权重、耦合或 MS-bar 因子。
因此这里比较的是同一个有限 BPH 条件，不能把裸 MS 极点直接充当有限振幅。

实际结果为
\begin{align*}
 c_0(0,1)&=-0.9987759652197846\ \pm3.25\times10^{-6},\\
 c_0((1,1,1),1)&=-1.0721646501342081\ \pm4.45\times10^{-6},\\
 J_R^{\rm native}((1,1,1),1)
 &=-0.07338868491442352\ \pm5.51\times10^{-6}.
\end{align*}
两次极点估计差为 $1.03\times10^{-6}$，其组合报告标准差为
$1.26\times10^{-6}$；极点重建的周期也与独立正生成树多项式采样一致。
通用图编译器使用 Hepp 采样、16 组各 $2^{22}$ 个点，得到
\[
 J_R^{\rm generic}=-0.07338791561208803\ \pm6.19\times10^{-6},
\]
半样本细化差为 $6.64\times10^{-6}$，两条方法的实际绝对差为
$7.69\times10^{-7}$。上述正负误差都是随机数值诊断，不是严格区间。

原始级数、种子、输入图、脚本、生成与编译日志、库哈希及时间均已保留。
执行期间的网页渲染修复改变过全包指纹；完整旧源码快照证实数学依赖逐文件未变，
比较记录同时保留旧执行指纹和当前指纹，没有覆盖原始变化标志。
联合 UV 算术修复后，上述通用 Hepp QMC 已在当前冻结源码下重新实际执行；
上游原始级数重解析不被称为又一次原生积分，旧通用 QMC 仍保留历史执行记录。

\subsection{真嵌套三圈图的完整积分试算}

另一次实际运行选择 \texttt{L3E3G53}，在目录边序中具有
$\{5,6\}\subsetneq\{0,1,3,5,6\}$ 两个严格嵌套的二点子图，
两者的减除阶数均为 $\omega=2$。
在 $M=1$、$(s_1,s_2,s_3)=(1,1,1)$ 下，
完整递归编译密度使用 8 组共 524288 个采样点，得到
\[
 J_R=-0.00023303604812\ \pm5.59\times10^{-5}.
\]
更换独立随机化及重要性参数，另用 8 组共 262144 个点，得到
$-0.00029981283515\ \pm2.52\times10^{-5}$。
两者差为 $6.68\times10^{-5}$，约为组合报告标准误的 1.09 倍；
两次半样本细化差分别为 $3.67\times10^{-5}$ 和 $1.91\times10^{-5}$。
这一结果证明程序已实际完成真嵌套图的完整积分试算，当前精度仍较粗。
两次计算复用同一图导出有限密度，属于独立随机化的一致性检查，
不是独立解析核互校，也没有给出严格积分误差界。

\subsection{四圈完整递归和的独立核对}

第\ref{sec:grouped-sampling}节的分层工具根据输入图识别三个不交bubble及收缩后的一圈结构，
完整密度仍由\texttt{compile\_bph}构造，不调用第\ref{sec:four-loop-independent}节的参考公式。
默认总尺度形状参数为1；经保留的pilot比较后，主核查明确使用总尺度形状$1/2$、
bubble内部Beta形状3及bare边比例形状1。
实际12组共196608点给出
\[
 J_{\rm generic}=-0.0005154702382511922\ \pm1.80923\times10^{-6}.
\]
独立外圈动量参考为$-0.0005158785412715456$，实际绝对差为$4.08303\times10^{-7}$；
主采样半样本细化差为$1.18853\times10^{-6}$。
独立种子8组共65536点给出
$-0.0005116648680529154\ \pm3.69832\times10^{-6}$，
两次随机化相差$3.80537\times10^{-6}$。
正负数值均指报告的随机组标准误，不是严格总误差区间。

普通Dirichlet及初始平面分组未达到精度的原始记录仍保留，
不会因最终提案达到目标而改写为成功。新的采样权重同时经过确定性单项式矩与独立PDF检查；
参考标量只在完成积分后比较。图号一致性修复后，主结果及独立种子已用正式工具重新实际执行，
工具起止SHA、完整输入图、程序SHA和冻结核心指纹均保存。
完整比较见\path{results/automatic_four_loop_independent.json}，
原始采样及失败对照见\path{docs/automatic-group-sampling-evidence/raw-files.json}。
该例提供一个实际四圈非primitive完整有限积分核对，并不代表所有四圈图均已精密求积。

\subsection{当前完成范围与未闭合部分}

现在已经有完整图编译、低圈完整有限积分互校、35 张高圈 primitive 解析证书，
以及一张三圈 primitive 的真正外部原生后端互校。
选定高圈实例的实际执行与完成状态单独记录。
其中 \texttt{L3E3G21} 与 \texttt{L4E3G205} 分别含两张、三张不交二点子图，
这两张图的完整积分检验的是多重不交二点反项，本身不构成真嵌套的证据；
真嵌套图的粗精度试算见上一小节，它仍缺少独立解析核互校。
四圈多重不交二点反项的完整独立核对见上一小节；选定高圈粗精度执行另见
\path{results/automatic_higher_integrals.json}。
点态 UV 检查不能替代整个参数积分控制，采样误差也不能替代原动量的绝对收敛证明。
698 张图的编译覆盖不等于 698 张图均已完成高精度全积分，
一般递归 UV 证明和任意圈归纳尚未闭合。

\subsection{复现入口与证据文件}

在解压后的软件目录运行。表达式构造和 primitive 证书只需 Python3.11 或更新版本；
通用积分使用可选 NumPy、SciPy，独立高精度微分参考需要 mpmath，
原生 Laurent 级数解析另需 pySecDec。
\begin{samepage}
\begin{verbatim}
python3 -m phi3lab build-automatic
python3 -m phi3lab automatic --loops 3 --legs 3 --index 5
python3 -m phi3lab automatic --input my-graph.json --integrate
python3 tools/compare_automatic_bph.py
python3 tools/verify_automatic_integrals.py
python3 tools/compare_automatic_native.py
python3 tools/run_automatic_higher.py
python3 tools/verify_automatic_four_loop.py
python3 tools/verify_tests.py
\end{verbatim}
\end{samepage}
自定义图 JSON 使用整数顶点字段 \texttt{vertices}、边列表 \texttt{edges}
和固定标签外腿列表 \texttt{external}，与目录中的图格式一致。
边的重排或顶点的重命名不会切换到按名字指定的专用积分公式。

\begin{itemize}
\item 图集与精确权重：\path{results/higher_graph_feyngen_comparison.json}。
\item 698 张图的表达式：\path{results/automatic/catalog.json} 及其逐图文件。
\item 44 项独立密度检查：\path{results/automatic_bph_comparison.json}。
\item 14 项完整两圈比较：\path{results/automatic_integral_comparison.json}。
\item 原生三圈有限转换与比较：\path{results/automatic_native_comparison.json}。
\item 真嵌套三圈完整积分试算：\path{results/automatic_nested_integral.json}。
\item 实际原生构建记录：\path{docs/automatic-native-build.json}。
\end{itemize}
每份证据的原始执行时间、源码指纹与实际工作阶段分别保存。
重新解析原始输出和重新执行积分是两个不同操作，不能混写。

\subsection{主要来源与方案边界}

历史 Henge 递归见 Kennedy、Binoth 与 Rippon，
\href{https://arxiv.org/abs/0712.1016}{arXiv:0712.1016}。
Brown 与 Kreimer 的
\href{https://arxiv.org/abs/1112.1180}{arXiv:1112.1180}
提供更广泛的参数重整化背景，但其 single-scale 构造没有被直接替换成本项目的完整有质量零动量条件。
现代 pySecDec 更新见 Heinrich 等，
\href{https://arxiv.org/abs/2305.19768}{arXiv:2305.19768}，
CPC 295 (2024) 108956；实际版本、输入和已执行点在各自记录中明确给出。
官方组合图参考为 \url{https://github.com/michibo/feyncop}。
Tropical 积分参考
\href{https://arxiv.org/abs/2302.08955}{arXiv:2302.08955}
要求合适的 quasi-finite 输入，不能据此直接把有子发散裸图当成已匹配 BPH 的参考。
本项目的递归编译器、Taylor jets 与采样代码由协作实现，没有打包第三方科学源码；
本文不宣称方法新颖性或一般圈数定理已经完成。

\end{document}
